\section{Marine Biofouling: Mechanisms and Engineering Consequences}
\subsection{Attachment sequence and fouling stages}
A useful engineering sequence is conditioning film, bacterial attachment, biofilm development, diatom and algal settlement, and macrofouler recruitment. The sequence is not deterministic: surface chemistry, roughness, flow, nutrient conditions, season, and organism physiology modify both timing and community composition. Comparative biofilm work reports different communities on biocidal and fouling-release coatings, demonstrating that control changes the assemblage rather than simply switching fouling on or off.

\subsection{Adhesion, wettability, roughness, and rheology}
Surface energy, hydration, chemical functionality, elastic response, nanoscale texture, and lubricant retention can alter attachment and removal. However, superhydrophobicity or a low contact angle is a property, not proof of antifouling performance. Biofilm rheology and its interaction with shear can matter as much as initial settlement \cite{snowdon_surface_2023}. Different organisms use different adhesion mechanisms, helping explain why results across diatoms, bacteria, algae, barnacles, and mixed fouling cannot be ranked on one scale.

\subsection{Roughness, drag, and powering}
Fouling and coating irregularities increase the hydrodynamic penalty through changes in near-wall flow and effective roughness. Grooming studies show lower drag for groomed than ungroomed biofilms, while operating-ship evidence links higher-quality coating and lower attachment with lower fuel consumption \cite{hearin_analysis_2016,hunsucker_biofilm_2018,hakim_investigation_2019}. These findings support the chain surface condition $\rightarrow$ fouling $\rightarrow$ roughness and adhesion $\rightarrow$ resistance $\rightarrow$ powering and fuel, but do not establish a universal conversion factor.

Translating a specific biological state to a precise drag penalty relies on the equivalent sand-grain roughness height ($k_s$), a hydrodynamic length scale rather than a direct geometric measurement, first conceptualized from rough-pipe and rough-plate flow \cite{nikuradse_1933,schlichting_1936}. Townsin established the ship-scale frictional-allowance framework later adopted by the ITTC, relating drag increments to the maximum peak-to-trough profile height over a 50~mm sampling length ($Rt_{50}$) \cite{townsin_mosaad_1985,townsin_1991,townsin_2003}; critically, Townsin noted that height-based formulations correlate reliably with drag only for moderately rough surfaces ($Rt_{50}<230\,\mu\text{m}$), a threshold that realistic biofouling routinely exceeds \cite{townsin_2003}. Granville's similarity-law scaling method provides the boundary-layer bridge from a measured roughness function ($\Delta U^+$) to full-scale resistance \cite{granville_1958,granville_1987}, and Candries demonstrated experimentally that friction on coated marine surfaces is not captured by vertical height alone: texture, slope, spatial frequency, and form each contribute independently, particularly when comparing self-polishing copolymers with fouling-release coatings \cite{candries_2001_thesis,candries_anderson_atlar_2001,candries_atlar_2005}. $k_s$ is commonly estimated from height statistics (for example $k_s\approx0.17\,Ra$) or from multi-parameter fits incorporating skewness \cite{schultz2007roughness,flack_schultz_2010}, but no universal correlation exists for the irregular, three-dimensional topographies produced by real fouling communities \cite{speranza_modelling_2019}. Biofilms and diatomaceous slimes are viscoelastic rather than rigid packed grains, so a static coupon measurement cannot be assumed to give the same $k_s$ as the same surface under turbulent shear that deforms or strips the fouling layer \cite{snowdon_surface_2023}.

Multiscale roughness and viscoelastic slime introduce an additional scale-transfer problem. A biofilm can alter near-wall streaks, effective mixing length, and apparent roughness non-monotonically as coverage and shear change. The evidence matrix contains no calibrated relation linking a particular LoF class or biofilm thickness to a ship-scale $k_s$ or $\Delta C_F$; those quantities are therefore NR. External turbulent-flow LIS measurements likewise indicate that theoretical slip benefits can be limited when turbulent Reynolds stresses strip or deform the lubricant layer \cite{arxiv170609934}. Coupon settlement reductions should not be converted linearly into ship drag without boundary-layer and resistance measurements.

\subsection{Operational dependence}
Static periods, speed, route, temperature, nutrient regime, niche geometry, and cleaning history govern exposure. Hull coatings, aquaculture nets, membranes, and heat exchangers impose different loading and access constraints. A valid comparison must therefore report asset, flow state, fouling community, coating age, cleaning history, and endpoint.

\section{Design Principles for Sustainable Control}
\subsection{Antifouling, fouling-release, anti-adhesion, and removal}
Antifouling seeks to prevent, inhibit, repel, or deter settlement; fouling-release weakens adhesion so motion or cleaning can remove organisms; anti-adhesion reduces attachment without necessarily killing; self-cleaning uses a stimulus or interface to remove deposits; mechanical management removes established fouling. Hybrid systems may combine modes, but the environmental and operational implications differ.

\subsection{The durability--performance trade-off}
Nanofillers can reinforce a polymer yet agglomerate at higher loading; lubricants can lower adhesion yet deplete or redistribute; hierarchical texture can improve wetting behavior yet be vulnerable to abrasion; and active chemistry can improve inhibition yet create release concerns. Cavas et al. show that carbon fillers improved PDMS mechanical properties without necessarily improving antifouling or cleaning performance \cite{cavas_reinforcement_2018}. The engineering objective is therefore retention of the relevant interface under service exposure, not the initial property alone.
